% !PRINTED 588
\weberpage{588}
%<<CONT
and in this way one determines a sequence of integers $\alpha, \alpha_1, \alpha_2, \alpha_3 \dots$
with steadily decreasing norm, which can be continued so long
as zero does not occur among them. Since these norms are positive
rational integers which always decrease, zero must appear
after a finite number of steps, and one thus obtains
a system of equations

\weq{1}{\begin{aligned}
\alpha &= \mu\alpha_1 + \alpha_2,\\
\alpha_1 &= \mu_1\alpha_2 + \alpha_3,\\
&\cdots\cdots\cdots\cdots\\
\alpha_{h-2} &= \mu_{h-2}\alpha_{h-1} + \alpha_h,\\
\alpha_{h-1} &= \mu_{h-1}\alpha_h.
\end{aligned}}

From this one concludes that the whole sequence of numbers
$\alpha, \alpha_1, \alpha_2 \dots$, and thus in particular also $\alpha$ and $\alpha_1$, are divisible
by $\alpha_h$, and conversely, that every common divisor of $\alpha$
and $\alpha_1$ is a divisor of all the following $\alpha$, and hence also of $\alpha_h$.
Every other number that has these two properties must
be associated with $\alpha_h$, and one therefore calls $\alpha_h$ (and every number
associated with $\alpha_h$) the \gesperrt{greatest common divisor} of
$\alpha$ and $\alpha_1$.

If from the algorithm (1) we eliminate the intermediate $\alpha$,
by substituting $\alpha_2$ from the first equation into the second, then $\alpha_3$
from the second into the third, there results
a result of the form $\alpha\varkappa + \alpha_1\lambda = \alpha_h$, and we can
therefore, if instead of $\alpha, \alpha_1, \alpha_h$ we write $\alpha, \beta, \delta$, state the theorem:

4. \gesperrt{If $\alpha, \beta$ are any two integers in $R(i)$,
and $\delta$ their greatest common divisor,
then the integers $\varkappa, \lambda$
can be so determined that}

\weq{2}{\alpha\varkappa + \beta\lambda = \delta}
\gesperrt{holds.}

If the greatest common divisor $\delta$ of two numbers
$\alpha, \beta$ is a unit, then $\alpha$ and $\beta$ are called relatively prime, and in this
case we can satisfy the equation

\weq{3}{\alpha\varkappa + \beta\lambda = 1}
by integral $\varkappa, \lambda$.

If $\beta$ is a prime, then either $\alpha$ is divisible by $\beta$,
or $\alpha$ and $\beta$ are relatively prime. In the latter case
equation (3) holds. If then $\gamma$ is another number in $R(i)$, it follows from (3) that
\[
\alpha\gamma\varkappa + \beta\gamma\lambda = \gamma.
\]
